\scanpage{44}\mspage{43}
\inserted{Cette relation est réflexive symétrique.} On vérifie \deleted{tout} de suite (en utilisant $M\Lambda\,1$ dans le cas de $Y\leq\begin{matrix}X\\[-.4ex]\rotatebox{90}{$\leq$}\\[-.4ex]X'\end{matrix}$) que \deleted{cette \unsure{relat}} \deleted{$\unsure{\mathcal R(X|\mathscr R(X}$} si $X'\leq X,\ Y'\leq Y$, \deleted{alors} $\mathcal R(X,Y)\Longrightarrow\mathcal R(X',Y')$ \deleted{\unsure{et l'}} \unsure{le'|l'}axiome $\unsure{ML|M\Lambda}\,5$ est \deleted{\unsure{de}} \unsure{satisfait|toujours} \unsure{tautologiquement|tautologique} pour cette condition. Ainsi l'axiome $M\Lambda\,5$ prend la forme
\[
X\mathrel{\rlap{$>$}<} Y\quad\Longrightarrow\quad\mathcal R(X,Y)
\]
et on voit donc que les relations cherchées $\mathrel{\rlap{$>$}<}$ sont les relations \inserted{réfl. sym.} ``saturées pour $\leq$''\unsure{)|,} qui impliquent $\mathcal R(X,Y)$. La relation admissible $\mathrel{\rlap{$>$}<}$ la plus large sera obtenue, en prenant $\mathcal R$ elle-\unsure{$=$|même}.

Dans le cas \inserted{d'un atlas} ensembliste spacieux, on a $\mathrel{\rlap{$>$}<}=\mathcal R$, mais pas dans d'autres cas intéressants. Ainsi on note que $\mathcal R(x,y)$ tjrs vrai si $x,y$ sont deux \unsure{$L$} --- or dans les cas non ensemblistes, il \unsure{pourra|arrivera} exister des lieux $x,y$ qui ne soient compatibles.

Je voudrais maintenant regarder le cas où $\mathrel{\rlap{$>$}<}$ peut se définir par la connaissance de $\parallel$ sur $\Lambda$, par une condition

\uline{$At\,C$}$_{\Lambda}$ (\unsure{espac|spac|espace}\deleted{\unsure{e}}) Soient $X,Y\in M$, \deleted{Alors $X\mathrel{\rlap{$>$}<} Y$}\inserted{$L=\widetilde X\cap\widetilde Y$.}
\[
\left|\begin{aligned}
X\mathrel{\rlap{$>$}<} Y\Longleftrightarrow{}&\deleted{\forall\,X',Y'\in\Lambda,\ X'\ll X,\ Y'\ll Y}\\
&\deleted{\forall\,X'\in \mathrm{omb}_{\Lambda}(X)\smallsetminus \mathrm{omb}_{\Lambda}(L)}\\
&\deleted{Y'\in \mathrm{omb}_{\Lambda}(Y)\smallsetminus \mathrm{omb}_{\Lambda}(L)}\\
&\inserted{\forall\,X'\in \mathrm{omb}_{\Lambda}X,\ Y'\in\mathrm{Omb}_{\Lambda}Y\ \text{avec }X'_L=Y'_L=\varnothing,}\\
&\text{on a}\quad X'\parallel Y'
\end{aligned}\right.
\]
(où on a posé, \deleted{pour $L\unsure{\subset|\in}M,\ \mathrm{omb}_{\Lambda}(L)=\bigcup_{Z\in L}\mathrm{omb}_{\Lambda}(Z)$} \inserted{$L=\widetilde X\cap\widetilde Y$, $X'_L=\deleted{\unsure{\widetilde X}}\widetilde{X'}\cap\mathrm{Omb}_{\deleted{\unsure{\Lambda}}}(L)=\left\{Z\,\middle|\,\begin{aligned}&Z\leq X'\\&\exists\,S\in L\\&\deleted{\unsure{\text{avec}}}Z\ll S\end{aligned}\right\}$}\unsure{\}}.

\uline{Cor} Soient $X,Y\in M$. \deleted{tels que} \deleted{Alors $X'$} \inserted{Alors} $X\parallel Y$ \deleted{i.e. $X\deleted{\unsure{\parallel}}\mathrel{\rlap{$>$}<} Y$ et $L=\varnothing$}. \deleted{Si $X'$} $\forall\, X'\in \mathrm{omb}_{\Lambda}(X)\quad\forall\,Y'\in \mathrm{omb}_{\Lambda}Y$ on a $X'\parallel Y'$.
